\section{Ioffe-time distributions and pseudo-PDFs}



 The basic object for defining parton distributions is a matrix element
 of a bilocal operator that  (skipping inessential details of  its
 spin structure)   may be written   generically  like
 $\langle  p | \phi (0)  \phi(z)  | p \rangle $.
Due to invariance under Lorentz transformations, it is given by a function of two scalars,
 the  {\it Ioffe time}   $(pz)$  \cite{Ioffe:1969kf}
 (which will  be denoted by $-\nu$)  and  the interval  $z^2$
  \begin{align}
  \langle p |   \phi(0) \phi (z)|p \rangle
=  & {\cal M} (-(pz), -z^2)  =  {\cal M} (\nu, -z^2)
\,
 \
\label{lorentz}
\end{align}
 (again, the sign for the second argument   is   chosen  so as to have
  a positive value  for \mbox{spacelike $z$)}.
 One can demonstrate \cite{Radyushkin:2016hsy,Radyushkin:1983wh}    that, for all relevant Feynman diagrams,
 its  Fourier transform  ${\cal P} (x, -z^2)$ with respect to $(pz)$
 has $-1 \leq x \leq 1$ as support, i.e.,
   \begin{align}
 {\cal M} (-(pz), -z^2)
&   =
 \int_{-1}^1 dx
 \, e^{-i x (pz) } \,  {\cal P} (x, -z^2)  \   .
  \label{MPD}
\end{align}
   In this  covariant  definition of  $x$,  one does not need to assume that
    $z$ is on the light cone
$z^2=0$ or that $p$   is light-like $p^2=0$.

  On the light cone $z^2=0$, we formally have
 \mbox{$   {\cal P} (x, 0) =f(x) $.}  Hence, the function $ {\cal P} (x, -z^2) $
 may be treated as   a generalization of the concept of
PDFs onto non-lightlike intervals $z^2$, and  following   \cite{Radyushkin:2017cyf} ,  we will  refer to it as the
 {\it pseudo-PDF}.  In view of lattice applications, we will take the separation
$z=\{0,0,0,z_3\}$ oriented  in the direction specified by the hadron momentum $p=\{E,0,0,P\}$.



  In renormalizable theories (including QCD),
  the function  $ {\cal M} (\nu, -z^2)  $
 has  logarithmic $\sim \ln (-z^2)$  singularities.
 In deep inelastic scattering (DIS), they  result
 in a logarithmic  scaling violation with respect
 to the photon virtuality  $Q^2$.
 A wide-spread  statement is that the \mbox{$Q^2$-dependent}   DIS structure functions
  $W(x_B,Q^2)$  probe the   hadron structure
 at distances $\sim 1/Q$.
 In the case  of the pseudo-PDFs ${\cal P} (x,z_3^2)$,
 one may say that  they literally describe the hadron structure
 at  the distance $z_3$.

 Just like the  DIS  form factors  $W(x_B,Q^2)$ are
 written in terms of  the universal parton densities $f(x, Q^2)$,
 the pseudo-PDFs   obtained   from
lattice calculations   may  be expressed through  the usual
 parton distributions. The  latter are  defined by
 the operators on the light cone $z^2=0$, i.e., in a
 logarithmically  singular limit.
 In the approach based on  the operator product expansion
(OPE),  the standard procedure is to  remove  these singularities  with the help of some
 prescription.

 The most popular of them is the  $\overline{\rm MS}$  scheme based
on the dimensional  regularization.
 Consequently,  the resulting PDFs  have a dependence on the renormalization scale
$\mu$, and therefore one should write the PDFs as $   f(x, \mu^2)$.
Switching from $x$  to the Ioffe time $\nu$ gives  the functions
 \begin{align}
 {\cal I} (\nu, \mu^2)
&   =
 \int_{-1}^1 dx
 \, e^{i x \nu  } \,  f (x, \mu^2)  \
  \label{ITD}
\end{align}
introduced in Ref.  \cite{Braun:1994jq}  and  called there  the {\it Ioffe-time distributions}.
In this context, the
functions  ${\cal M}(\nu , -z^2) $   that are  the  Fourier transforms of
pseudo-PDFs,    should be called the  Ioffe-time {\it pseudo-distributions}
or {\it pseudo-ITDs}.



To get  a relation between the pseudo-PDFs ${\cal P} (x,z_3^2)$  and the
$\overline{\rm MS}$     parton densities $ f(x, \mu^2)$, one can use
 the  nonlocal  light-cone OPE \cite{Anikin:1978tj,Balitsky:1987bk}
 (see also \cite{Radyushkin:2017lvu})
for  the  matrix element defining
${\cal P} (x,z_3^2)$, i.e., for  the pseudo-ITD. The result
\begin{align}
{\cal M } (\nu, -z^2 ) =&  \sum_i  \int_{-1}^1 dw \, C_i (w, z^2 \mu ^2, \alpha_s) \,  {\cal I}_i (w \nu,\mu^2)
\nn & + {\cal O} (z^2)
  \  ,
\label{OPE}
\end{align}
has the   structure  similar to that of the usual  OPE  for the DIS structure functions $W (x, Q^2)$.
In this expression,   the twist-2 coefficient functions $C_i$ are
given by an expansion in the strong coupling constant $\alpha_s$,
while  ${\cal O} (z^2) $ symbolizes  higher-twist terms.

However, the application of the OPE
to the pseudo-ITDs  and pseudo-PDFs   in QCD  faces
complications   related to the gauge link.  Namely, when $z$ is off the light cone,
the link  generates
linear $\sim z_3/a$ and logarithmic $\sim \ln (1+z_3^2/a^2)$  ultraviolet (UV) divergences,
where $a$ is an UV regulator with the dimension of length
(it may be a finite lattice  spacing).
Though
 disappearing  for $z_3=0$, these   divergences
require an additional UV  regularization when $z_3$ is finite.


Fortunately, these divergences are multiplicative \mbox{\cite{Polyakov:1980ca,Dotsenko:1979wb,Brandt:1981kf,Aoyama:1981ev,Craigie:1980qs}}  (see also
 recent Refs.
 \cite{Ishikawa:2017faj,Ji:2017oey,Green:2017xeu}), and cancel in the ratio,
 the  reduced Ioffe-time distribution,
 \begin{align}
{\mathfrak M} (\nu, z_3^2) \equiv \frac{ {\cal M} (\nu, z_3^2)}{{\cal M} (0, z_3^2)} \   ,
 \label{redm}
\end{align}
 introduced in our paper \cite{Radyushkin:2017cyf}, and partially motivated
 by this cancellation.
 The remaining $\ln z_3^2$ singularities, present only in the numerator of the ratio,
  are   described by the nonlocal light-cone OPE.


As stated in Ref.  \cite{Orginos:2017kos}, for
 small spacelike intervals $z^2=-z_3^2$,  and at the leading logarithm level, the reduced pseudo-PDFs are
related to the $\overline{\rm MS}$ distributions by a simple rescaling of their second arguments,
namely,
\begin{align}
\mu^2  =4e^{-2\gamma_E}/ z_3^2
  \  ,
\label{Ptof}
\end{align}
where $\gamma_E$  is the Euler's constant (a more detailed discussion  will be given later on).
This  rescaling  factor  is very close to 1, since \mbox{$2e^{-\gamma_E}=1.12$. }
However, this factor may be changed by the $ {\cal O} (\alpha_s) $  terms present in the coefficient function.
